\chapter{Introduction}
\label{ch:introduction}

\section{Motivation and Objectives}
\label{sec:intro:motivation}

The problem addressed by this thesis arises at the intersection of two
well-documented shortcomings in current practice. On one side, conventional
interior design tools, whether two-dimensional floor plans, printed catalogues,
or digital product configurators, present furniture and spatial layouts as
abstractions that strip away the perceptual properties that ultimately determine
whether a redesign succeeds: how an object occludes a window, whether a sofa
visually dominates a narrow corridor, or how a chosen material interacts with
the room's natural light. On the other side, consumer augmented reality
applications that \emph{do} allow virtual objects to be previewed in a live
camera feed share a structural limitation identified empirically by
Kandil~et~al.~\cite{kandil2021ux}: they treat the room as a transiently
detected floor plane, provide no persistent room model, and require the user
to be physically present in the space for every design iteration. The result
is a tool that is useful for a single in-store moment but not for the extended,
iterative process that characterises real interior redesign decisions.

A more principled approach is to treat the room scan itself as the primary
artefact of the workflow. If the real room is captured once as a metric,
semantically annotated three-dimensional model, a \emph{digital twin}, subsequent
design iterations can be performed at any time and from any location without
requiring the user to return to the physical space. Indoor scene understanding
research has demonstrated that this level of structure is achievable at scale
on consumer hardware: ScanNet~\cite{dai2017scannet} showed that richly annotated
RGB-D reconstructions with per-instance semantic labels support a wide range of
downstream tasks, and Matterport3D~\cite{CAS+17} extended that work to
building-scale scenes with 40-category instance labels. Apple's RoomPlan
API~\cite{apple2024roomplan2} brings a real-time, on-device subset of this
capability to supported iPhones and iPads, outputting a parametric
\texttt{CapturedRoom} whose typed structural surfaces and classified furniture
instances reflect the same conceptual taxonomy that these research datasets
made trainable.

A second motivation concerns the asset authoring bottleneck. Replacing a
detected furniture item with a virtual equivalent requires a three-dimensional
model of that item. Static libraries cover a limited catalogue and exclude
bespoke or non-commercial pieces. Object Capture~\cite{apple2024roomplan2}
enables the user to scan a physical object with the same device and obtain a
USDZ mesh through an on-device photogrammetric pipeline, closing the loop
between room capture and asset creation without external tooling during the
core staging workflow.

A third motivation concerns representation choice and workflow completeness.
3D Gaussian Splatting~\cite{kerbl20233dgs} achieves real-time rendering quality
competitive with neural radiance fields while remaining differentiable for
offline training. In this application, explicit meshes remain the medium for
interactive furniture replacement in RealityKit, whereas imported splats and
feedforward SHARP reconstructions serve as \emph{visual benchmarks}. Users who
require desktop refinement can export staged rooms to Blender as glTF-based ZIP
archives; users who need rapid radiance-field previews can invoke a local SHARP
inference server. Chapter~\ref{ch:theory2} formalises the three object-level
reconstruction paths (Assets~A, B, and~C); Chapter~\ref{ch:implementation}
documents their implementation.

The work pursues the following objectives, each directly justified by the
gaps identified above:

\begin{enumerate}
  \item \textbf{Room-scale semantic capture.}
    Integrate RoomPlan so that a scan yields a \texttt{CapturedRoom} containing
    typed structural elements (walls, floors, ceilings, doors, windows) and
    classified furniture instances with metric transforms and bounding dimensions.
    \emph{Justification:} without a persistent, semantically structured room
    model, furniture replacement cannot be grounded to real geometry.

  \item \textbf{Persistent virtual staging.}
    Provide a non-AR three-dimensional editor with orbit-style navigation and
    gesture-based manipulation (translation, rotation, scale), with state
    persisted across application launches.
    \emph{Justification:} Kandil~et~al.~\cite{kandil2021ux} found that
    requiring physical presence and absence of session persistence were primary
    usability barriers in existing AR interior apps.

  \item \textbf{Object-scale asset authoring (Asset~A).}
    Integrate Object Capture for guided multi-view photogrammetric
    reconstruction into USDZ, extending the furniture library beyond static
    catalogue content.

  \item \textbf{Radiance-field benchmarks (Assets~B and~C).}
    Import offline-trained 3DGS \texttt{.ply} files (Asset~B) and invoke Apple's
    SHARP feedforward model on a LAN server from a single photograph (Asset~C);
    render both through a shared Metal splat viewer for qualitative comparison
    with meshes---not for furniture editing.

  \item \textbf{Room analytics and planning aids.}
    Compute metric dimensions and renovation estimates from scan geometry;
    provide rule-based interior-design suggestions (including lighting density
    heuristics) in a dedicated statistics interface.

  \item \textbf{DCC hand-off.}
    Export completed rooms as glTF-based ZIP archives with USDZ furniture,
    manifests, and a Blender helper script for further editing outside the device.
\end{enumerate}


\section{Contributions and Scope}
\label{sec:intro:contributions}

The contributions of this thesis are primarily integrative and
engineering-oriented, consistent with the scope of a bachelor project:

\begin{itemize}
  \item \textbf{A unified scan-to-edit pipeline.}
    A complete iOS application connecting RoomPlan capture to a serialised
    room data model, to a RealityKit scene graph, and to furniture replacement
    with persistence across application launches.

  \item \textbf{Theoretical framing of Apple's capture APIs within the
    academic literature.}
    A structured account of how Object Capture instantiates the classical
    photogrammetric chain and how RoomPlan builds on depth fusion in the TSDF
    family~\cite{newcombe2011kinectfusion}, augmented by learned structure and
    semantics, with a principled comparison to the 3D Gaussian Splatting
    rendering path~\cite{kerbl20233dgs}.

  \item \textbf{A three-way comparative analysis} of mesh (Asset~A),
    offline 3DGS (Asset~B), and SHARP splats (Asset~C), clarifying roles:
    interactive staging versus visual benchmarking.

  \item \textbf{A Metal-based Gaussian splat renderer for iOS.}
    Loading standard binary PLY files from common training frameworks at
    interactive frame rates, bridged into SwiftUI via \texttt{MTKView}.

  \item \textbf{Supporting services.}
    Two-level \texttt{CapturedRoom} caching, rule-based
    \texttt{SmartSuggestionsService}, scan-time colour sampling, Blender export,
    and local-network SHARP integration without third-party cloud upload for
    Apple capture paths.
\end{itemize}

The thesis focuses on software architecture, core algorithms at a conceptual
level, and the iOS platform. It does not contribute novel learning
architectures for room understanding: RoomPlan's neural models are Apple's
proprietary intellectual property. 3DGS training is assumed to run off-device
in the current system; on-device splat training remains
research-grade~\cite{Guo26}. Full LiDAR-quality RoomPlan output requires
iPhone~12~Pro or later, or a LiDAR-equipped iPad~Pro; the application
degrades gracefully on non-LiDAR devices where Apple's API permits.

\section{Document Structure}

Chapter~\ref{ch:theory1} presents photogrammetry, depth fusion, and the mapping to Object Capture and RoomPlan, followed by related work. Chapter~\ref{ch:theory2} introduces 3D Gaussian Splatting, feedforward SHARP inference, the mesh-versus-splat design choice, and the three-asset comparison. Chapter~\ref{ch:requirements} formalises functional and non-functional requirements and system constraints. Chapter~\ref{ch:implementation} describes architecture, data flow, implementation details, analytics, Blender export, the user interface, and testing. Chapter~\ref{ch:conclusions} summarises results, limitations, and future work.

The opening \textit{Context and Related Work} section situates the project against ScanNet-style indoor semantics~\cite{dai2017scannet}, AR interior design systems~\cite{kan2021ar}, and mobile AR usability evaluation~\cite{kandil2021ux}.
